\documentclass[11pt,a4paper]{article}
\usepackage[margin=25mm]{geometry}
\usepackage{amsmath,amssymb,amsthm,booktabs,microtype}
\usepackage[colorlinks=true,linkcolor=blue!50!black,urlcolor=blue!50!black]{hyperref}
\usepackage[dvipsnames]{xcolor}
\usepackage{array}
\usepackage{fancyhdr}
\pagestyle{fancy}
\fancyhf{}
\fancyhead[L]{\small Newton's Cathedral --- mathematical working note}
\fancyhead[R]{\small 8 November 2026}
\fancyfoot[C]{\small \thepage}
\renewcommand{\headrulewidth}{0.35pt}
\newtheorem{theorem}{Theorem}
\newtheorem{corollary}{Corollary}
\newtheorem{remark}{Remark}
\title{\vspace{-1.1cm}\textbf{Co-closed Curvature Coercivity and\\Holonomy-Stable 24-Track Transport}\\[5pt]\large A five-Platonic Hodge theorem, with an independent covariant-clock extension}
\author{Newton's Cathedral --- independent research note}
\date{8 November 2026 \;\textbullet\; Working draft, not a continuum Yang--Mills theorem}
\begin{document}
\maketitle
\vspace{-7mm}
\begin{abstract}
For the direct sum of the standard, unit-weight oriented cellular cochain complexes of all five Platonic solids, viewed as cellulations of $S^2$, we prove a sharp spectral inequality for co-closed edge cochains $\phi$. With $d_0:C^0\to C^1$ and $d_1:C^1\to C^2$, the inequality is
\[
\boxed{\quad d_0^{\!*}\phi=0\quad\Longrightarrow\quad
\|d_1\phi\|^2\ge (3-\sqrt 5)\,\|\phi\|^2.\quad}
\]
The coefficient is the exact smallest positive eigenvalue of the dodecahedral graph Laplacian, which is dual to the icosahedral face Laplacian. We verify the exact incidence spectrum, extend the bound to 24 arbitrary orthogonal clock transports, and compute the golden-angle helical dispersion. This concerns fixed finite complexes and does not establish scale-uniform lattice Yang--Mills coercivity, non-Abelian gauge dynamics, or a continuum mass gap.
\end{abstract}

\section*{Acknowledgment and attribution}
This independent Newton's Cathedral note was motivated by a forest-gauge
block Poincar\'e inequality shared from the Yang--Mills programme of
\textbf{James Montgomery Lockwood}. Lockwood's documented contributions
include constrained-cube analysis and its code and Lean~4 formalisations,
and the composition of metrics across scales. The associated paper is
jointly credited to \textbf{Deep Bhattacharjee, Ushashi Bhattacharya,
James Montgomery Lockwood and Tristan Romanov}~\cite{bhattacharjee2026}.
The specific $537/5000$ forest-gauge constant is quoted as supplied, not
independently established here. The standard Hodge and finite spectral
calculation that follows is independently defined and does not borrow
those authors' proofs or code. This research note is not coauthored or
endorsed by the Yang--Mills paper's contributors.

\section{Question and precise scope}
The shared Lockwood constrained-block claim is the inequality
\begin{equation}\label{eq:lockwood}
 \|d\varphi\|^2\ge \frac{537}{5000}\|\varphi\|^2
 \qquad\left(\frac{537}{5000}=0.1074\right)
\end{equation}
for forest-gauge link fields satisfying specific zero block-average constraints, uniformly over specified finite supports. The supplied image does not define those supports, gauge-fixing operators, or block constraints. Therefore its constant cannot be independently certified from the image alone. Lockwood's constrained-cube work is credited in the 2026 lattice Yang--Mills preprint of Bhattacharjee, Bhattacharya, Lockwood and Romanov~\cite{bhattacharjee2026}. The preprint itself distinguishes its conditional continuum criterion from a completed Yang--Mills solution.

Rather than copy equation~\eqref{eq:lockwood} into another system, we test the corresponding coercivity question on a fully specified independent object: the complete five-solid Platonic Hodge complex of Newton's Cathedral, of dimensions $(50,90,50)$ in cochain degrees $0,1,2$.

\section{Five-solid complex and exact theorem}
For each Platonic solid $P$ use the canonical oriented face-edge-vertex incidences, unit Euclidean inner products in cellular bases, and coboundaries
\[
 C^0(P)\xrightarrow{d_0}C^1(P)\xrightarrow{d_1}C^2(P),
 \qquad d_1d_0=0.
\]
The total complex is the direct sum of the tetrahedral, cubical, octahedral, dodecahedral and icosahedral cellulations. It contains $50$ vertices, $90$ edges and $50$ faces. The edge space splits orthogonally into $45$ exact and $45$ coexact dimensions.

\begin{theorem}[Sharp five-Platonic co-closed curvature bound]\label{thm:gap}
Let $\phi\in C^1$ for the full direct sum, and impose $d_0^{\!*}\phi=0$. Then
\begin{equation}\label{eq:theorem}
 \langle \phi,d_1^{\!*}d_1\phi\rangle
 =\|d_1\phi\|^2\ge \mu_*\|\phi\|^2,
 \qquad \mu_*:=3-\sqrt5=0.7639320225\ldots .
\end{equation}
The constant $\mu_*$ is optimal.
\end{theorem}

\begin{proof}
Each boundary surface is a cellular $2$-sphere. Hence its first cohomology vanishes, and
$\ker d_1=\operatorname{im}d_0$ and $C^1=\operatorname{im}d_0\oplus\operatorname{im}d_1^{\!*}$ orthogonally. The co-closed condition $d_0^{\!*}\phi=0$ therefore restricts $\phi$ to $\operatorname{im}d_1^{\!*}$. On that subspace the positive spectrum of $d_1^{\!*}d_1$ equals that of $d_1d_1^{\!*}$. Because every edge has exactly two incident faces with opposite boundary signs, $d_1d_1^{\!*}$ is exactly the combinatorial graph Laplacian of the \emph{dual} Platonic graph.

The dual-graph Laplacian gaps are listed in Table~\ref{tab:gaps}; their minimum is $3-\sqrt5$, occurring for the dodecahedral dual graph of the icosahedral cellulation. Exactness follows directly from the dodecahedral Laplacian characteristic polynomial
\begin{equation}\label{eq:charpoly}
 \det(\lambda I-L_{\mathrm{dodec}})
 =\lambda(\lambda-5)^4(\lambda-3)^4(\lambda-2)^5
 (\lambda^2-6\lambda+4)^3.
\end{equation}
Its least nonzero root is $3-\sqrt5$ (multiplicity three). Taking an associated coexact singular vector on the icosahedral summand and zero on the remaining four summands saturates~\eqref{eq:theorem}.
\end{proof}

\begin{table}[ht]
\centering\small
\begin{tabular}{lrrrl}
\toprule
Cell surface & $V$ & $E$ & $F$ & Smallest positive eigenvalue of $d_1^{\!*}d_1$ \\
\midrule
Tetrahedron &4&6&4& $4$ \\
Cube &8&12&6& $4$ \\
Octahedron &6&12&8& $2$ \\
Dodecahedron &20&30&12& $5-\sqrt5=2.763932\ldots$ \\
Icosahedron &12&30&20& $3-\sqrt5=0.763932\ldots$ \\
\midrule
Total & \textbf{50}&\textbf{90}&\textbf{50}& \textbf{Minimum $3-\sqrt5$} \\
\bottomrule
\end{tabular}
\caption{Exactly computed unit-weight incidence spectra. These are different operators and constraints from Lockwood's forest-gauge block problem.}\label{tab:gaps}
\end{table}

\section{A mathematically valid transfer consequence}
Set $K=d_1^{\!*}d_1$ restricted to $\ker d_0^{\!*}$ and take the \emph{explicitly defined} heat transfer $T_t=\exp(-tK)$ for $t>0$. Theorem~\ref{thm:gap} immediately gives
\begin{equation}\label{eq:semigroup}
 \|T_t\|_{\ker d_0^{\!*}}\le e^{-t(3-\sqrt5)},
 \qquad -\log T_t=tK\succeq t(3-\sqrt5)I.
\end{equation}
This supplies a finite-cell contraction estimate \emph{if} Cathedral transfer is defined from this $K$ or proved to obey a compatible operator inequality. It does not by itself identify $T_t$ with a previously proposed positive URT transfer $B^{\!*}B$ or with a physical time-evolution operator.

\begin{corollary}[Uncoupled 24-track replication]
For $24$ independent copies, $K_{24}=I_{24}\otimes K$ has the same sharp lower bound $3-\sqrt5$. Adding any positive semidefinite cross-track quadratic form preserves the bound. Indefinite couplings require a separate lower bound; commutativity is not necessary for a positive-semidefinite addition.
\end{corollary}


\section{Covariant 24-track transfer: a transport-stability theorem}
The full co-closed edge carrier $\mathcal H=\ker(d_0^{\!*})$ has dimension $45$.
Write $K=d_1^{\!*}d_1\big|_{\mathcal H}$, and let
$\boldsymbol\phi=(\phi_0,\ldots,\phi_{23})\in\mathcal H^{24}$.
Choose arbitrary orthogonal maps $U_j:\mathcal H\to\mathcal H$ with
$U_j^{\!*}U_j=I$, and a nonnegative coefficient $\kappa\ge0$. Define
\begin{equation}\label{eq:covenergy}
\mathcal E_{\kappa,U}(\boldsymbol\phi)=
\sum_{j=0}^{23}\|d_1\phi_j\|^2
+\kappa\sum_{j=0}^{23}\|\phi_{j+1}-U_j\phi_j\|^2,
\qquad \phi_{24}:=\phi_0.
\end{equation}
The maps $U_j$ are mathematical parallel transports on the co-closed carrier;
no physical gauge connection has been inferred from the finite incidence matrices.

\begin{theorem}[Covariant-clock coercivity, arbitrary orthogonal holonomy]
\label{thm:clock}
For every choice of the $24$ orthogonal maps, every $\kappa\ge0$, and
all $\boldsymbol\phi\in\mathcal H^{24}$,
\begin{equation}\label{eq:clockgap}
\boxed{\quad \mathcal E_{\kappa,U}(\boldsymbol\phi)
\ge (3-\sqrt5)\sum_{j=0}^{23}\|\phi_j\|^2.\quad}
\end{equation}
The constant is best possible \emph{uniformly over} the allowed transports.
Moreover, equality holds for a nonzero field precisely when each $\phi_j$
lies in the minimum-eigenvalue subspace of $K$ and every covariant link
vanishes: $\phi_{j+1}=U_j\phi_j$.
\end{theorem}
\begin{proof}
Apply Theorem~\ref{thm:gap} to each $\phi_j$ and sum. The $24$
additional terms are nonnegative. If equality holds, both sets of
nonnegative deficits must vanish, proving the stated criterion.
For $U_j=I$ and a constant sequence $\phi_j=\phi_*$, where
$K\phi_*=(3-\sqrt5)\phi_*$, equality is attained. Therefore no larger
constant can hold for \emph{all} orthogonal transports.
\end{proof}

This bound does not depend on commutativity of the $U_j$ or the choice of
clock holonomy $U_{23}\cdots U_0$. In particular, it is not obtained by
assuming that the original microscopic URT operator equals this quadratic
clock action; that identification must be tested separately.

\begin{corollary}[Exact golden-phase helical spectrum in a two-mode sector]
\label{cor:golden}
Take two orthonormal eigenvectors $a,b$ in the three-dimensional
$\mu_*$-eigenspace of $K$, with $\mu_*=3-\sqrt5$.
Let all $U_j=U$ rotate $\operatorname{span}\{a,b\}$ by the frozen
angle $\theta=2\pi/\phi^2$ and act as identity on its orthogonal
complement. On the complexified two-plane, the clock energy operator has
exact eigenvalues
\begin{equation}\label{eq:clockdisp}
\lambda_{m,\pm}=\mu_*+
4\kappa\sin^2\!\left(\frac{2\pi m/24\mp\theta}{2}\right),
\qquad m=0,\ldots,23.
\end{equation}
With $\kappa=1$, the smallest such helical eigenvalue is at $m=9$
(or its complex-conjugate partner),
\[
\lambda_{\mathrm{helix,min}}=0.765847419253977\ldots .
\]
The third, untouched $\mu_*$ eigenvector tensored with a constant
clock state still saturates the \emph{full} bound~\eqref{eq:clockgap};
therefore $0.765847\ldots$ is a \emph{two-plane sector} value, not
an improved global spectral gap.
\end{corollary}
\begin{proof}
In the two-plane let $U$ have eigenvalues $e^{\pm i\theta}$.
The periodic shift has eigenvalues $e^{2\pi im/24}$.
The covariant clock Laplacian is $(S-U)^{\!*}(S-U)$ on that sector;
its eigenvalues are
$|e^{2\pi im/24}-e^{\pm i\theta}|^2$, which yields
\eqref{eq:clockdisp}. The evaluated minimum is verified in the
accompanying executable code.
\end{proof}

\begin{corollary}[An explicit obstruction test for signed interactions]
\label{cor:signed}
Let $A$ denote the operator of~\eqref{eq:covenergy}
restricted to $\mathcal H^{24}$, and let $W=W^{\!*}$ be any additional
finite interaction (for example one proposed from the $H_{13}\oplus
H_{51}$ history sector). If $W_-$ is the positive negative-part
operator of $W$, then
\[
\lambda_{\min}(A+W)\ge \mu_*-\|W_-\|_{\rm op}.
\]
Hence $\|W_-\|_{\rm op}<3-\sqrt5$ suffices for strict coercivity.
If a hidden sector has $D\succeq\nu I>0$ and is coupled by $C$,
the eliminated Schur operator obeys the testable sufficient bound
\[
A-CD^{-1}C^{\!*}\succeq
\left(\mu_*-\frac{\|C\|_{\rm op}^2}{\nu}\right)I.
\]
These are \emph{conditional algebraic tests}, not estimates of actual
source-loaded Cathedral couplings: the operators $C,D$ and normalization
must be independently identified before numbers can be inserted.
\end{corollary}

\section{What the comparison does and does not establish}
The numerical ratio $(3-\sqrt5)/(537/5000)\approx7.11296$ is \emph{not} an improvement over Lockwood's inequality: it compares distinct configurations, normalizations and constraint spaces. Also, the un-gauged version of our inequality is false. For every nonconstant vertex function $f$, the nonzero edge field $\phi=d_0f$ has $d_1\phi=d_1d_0f=0$. The co-closed restriction is essential, as forest-gauge and block constraints are essential in the quoted research.

Three ingredients remain separate from this finite-dimensional result: (i) an explicit map from Cathedral carrier degrees of freedom (including $H_{13}\oplus H_{51}$ and the $24$-track ordering) to a specified lattice gauge theory; (ii) a coupling-dependent, scale-uniform coercivity or correlation estimate through refinement; and (iii) construction and nontriviality of the requisite four-dimensional continuum quantum theory with reflection positivity and a physical spectral gap. None is proved here. The boundedness of a single fixed finite Hodge operator must not be mistaken for uniformity as the lattice spacing tends to zero.

\section{Reproducibility and provenance}
The checks are implemented in two standalone scripts:
\begin{quote}\small\ttfamily
cathedral\_hodge\_gap.py\\
cathedral\_24track\_covariant\_gap\_verify.py
\end{quote}
Both independently construct the oriented boundary/coboundary matrices from each graph's planar embedding. They check $d_1d_0=0$, Euler characteristic, the ranks $\operatorname{rank}d_0=V-1$ and $\operatorname{rank}d_1=F-1$, compute the five spectra and symbolically factor~\eqref{eq:charpoly}. The second script checks the 24-track covariant energy for a golden-angle two-plane and 24 independently randomized orthogonal link transports. No external Yang--Mills proof or code is incorporated. The calculation uses only standard graph and cellular Hodge theory; novelty relative to existing mathematical literature is not claimed.

\begin{thebibliography}{9}
\bibitem{bhattacharjee2026} Deep Bhattacharjee, Ushashi Bhattacharya, \textbf{James Montgomery Lockwood} and Tristan Romanov, \emph{One Cube Estimate Implies the Yang--Mills Mass Gap}, Hexagon preprint 2610.00341v7 (2026), \url{https://hexagonmath.org/2610.00341v7}. The paper states that the condition required throughout continuum tuning and the construction of a nontrivial continuum limit remain open.
\end{thebibliography}
\end{document}
